%!TEX program = xelatex
\PassOptionsToPackage{full}{textcomp}
\documentclass[lang=en,11pt,a4paper,citestyle=authoryear,nobib=true]{elegantpaper}

\title{HW9 - STAT 709}
\author{Wells Guan}

\usepackage{array,url,stix}
\usepackage{tikz}
\usepackage{tikz-cd}
\usepackage{enumitem}
\usepackage{amsmath}
\usepackage{amssymb}
\usepackage{listings}
\usepackage{subfigure}

\lstset{
    frame=single,
    columns=fullflexible,
    keepspaces=true,
    showstringspaces=false,
    breaklines=true,
    breakatwhitespace=false,
    basicstyle=\ttfamily\footnotesize,
    numbers=left,
    numberstyle=\tiny,
    numbersep=6pt
}

\newcommand{\Z}{\mathbb{Z}}
\newcommand{\R}{\mathbb{R}}
\newcommand{\N}{\mathbb{N}}
\newcommand{\B}{\mathcal{B}}
\newcommand{\E}{\mathbb{E}}
\renewcommand{\P}{\mathbb{P}}
\newcommand{\norm}[1]{\left\lVert#1\right\rVert}

\begin{document}


\maketitle

\noindent Complete the required reading and all required exercises below.
Only Exercise 2.38 is optional.

\subsection*{Problem 1. Required reading}
Vershynin, \textit{High-Dimensional Probability}, second edition
(online PDF dated October 23, 2026): Sections 2.3 and 2.6--2.9.
All exercise numbers below refer to this edition.

\par\vspace{0.5em}\noindent\textbf{Reading notes.}\par
For independent Bernoulli random variables $X_i$, denote
$S_N=\sum_{i=1}^N X_i$ and $\mu=\E S_N$. We know the Chernoff bound
\[
\P(S_N\ge t)\le e^{-\mu}\left(\frac{e\mu}{t}\right)^t,
\qquad t\ge\mu>0.
\]
The same expression bounds $\P(S_N\le t)$ when $0<t\le\mu$.

For a real-valued random variable $Z$, write
\[
\norm{Z}_{\psi_2}
=\inf\{s>0:\E e^{Z^2/s^2}\le2\}.
\]
A finite $\psi_2$ norm is equivalent, up to absolute constants, to a
Gaussian tail bound and to $\norm{Z}_q\le C\norm{Z}_{\psi_2}\sqrt q$
for every $q\ge1$. If $\E Z=0$, it also gives
\[
\E e^{\lambda Z}\le
\exp\bigl(C\lambda^2\norm{Z}_{\psi_2}^2\bigr),
\qquad \lambda\in\R.
\]
In particular, if $Z_1,\ldots,Z_N$ are independent, mean-zero and
subgaussian, then
\[
\P\left(\left|\sum_{i=1}^N a_iZ_i\right|\ge t\right)
\le2\exp\left(-\frac{ct^2}
{\sum_{i=1}^N a_i^2\norm{Z_i}_{\psi_2}^2}\right).
\]

For subexponential random variables, the corresponding norm is
\[
\norm{Z}_{\psi_1}=\inf\{s>0:\E e^{|Z|/s}\le2\}.
\]
Notice that $\norm{Z^2}_{\psi_1}=\norm{Z}_{\psi_2}^2$.
For independent, mean-zero subexponential $Z_i$, Bernstein inequality gives
\[
\P\left(\left|\sum_{i=1}^N Z_i\right|\ge t\right)
\le2\exp\left[
-c\min\left\{\frac{t^2}{\sum_{i=1}^N\norm{Z_i}_{\psi_1}^2},
\frac{t}{\max_i\norm{Z_i}_{\psi_1}}\right\}\right].
\]
Here $c,C>0$ are absolute constants. In the two tail bounds above,
the case of a zero denominator is treated separately: the corresponding
sum is zero almost surely.

\subsection*{Problem 2. Concentration inequalities and Gaussian maxima}

\subsubsection*{Exercise 2.21. Boosting randomized algorithms (required)}
Imagine we have an algorithm for solving some decision problem, like
checking if $p$ is prime. Suppose the algorithm makes decisions randomly,
giving the correct result with probability $\frac12+\epsilon$ for some
$\epsilon>0$, only slightly better than pure guessing. To improve the
performance, run the algorithm $N$ times and take the majority vote.
Show that, for any $\delta\in(0,1)$, the answer is correct with probability
at least $1-\delta$, as long as
\[
N\ge\frac{1}{2\epsilon^2}\ln\left(\frac1\delta\right).
\]

\par\vspace{0.5em}\noindent\textbf{Sol.}\par
For a fixed input, run the algorithm with independent randomness and
denote by $X_i$ the indicator that the $i$th answer is correct. We know
that $X_i$ are independent Bernoulli random variables with mean
$\frac12+\epsilon$. Denote $S_N=\sum_{i=1}^N X_i$, then
\[
\E S_N=N\left(\frac12+\epsilon\right).
\]
If $S_N>N/2$, the majority answer is correct. Counting a tie as an error,
we have by the one-sided Hoeffding inequality
\[
\begin{aligned}
\P(\text{the majority answer is incorrect})
&\le\P(S_N\le N/2)\\
&=\P(S_N-\E S_N\le-N\epsilon)\\
&\le\exp(-2N\epsilon^2).
\end{aligned}
\]
For $N\ge(2\epsilon^2)^{-1}\ln(1/\delta)$, the last expression is at most
$\delta$, and hence the answer is correct with probability at least
$1-\delta$.
\vspace{1em}

\subsubsection*{Exercise 2.14. Chernoff inequality: small deviations (required)}
Let $X_i$ be independent Bernoulli random variables with parameters $p_i$.
Consider their sum $S_N=\sum_{i=1}^N X_i$ and denote its mean by
$\mu=\E S_N$. Show that
\[
\P(|S_N-\mu|\ge\delta\mu)
\le2\exp\left(-\frac{\delta^2\mu}{2+\delta}\right)
\qquad\text{for every }\delta\ge0.
\]

\par\vspace{0.5em}\noindent\textbf{Sol.}\par
If $\mu=0$ or $\delta=0$, the right side is $2$ and the result is
immediate. Assume $\mu>0$ and $\delta>0$. By independence,
\[
\begin{aligned}
\E e^{\lambda S_N}
&=\prod_{i=1}^N\bigl(1+p_i(e^\lambda-1)\bigr)\\
&\le\exp\bigl(\mu(e^\lambda-1)\bigr),
\qquad\lambda\in\R.
\end{aligned}
\]
Taking $\lambda=\ln(1+\delta)>0$ in Markov's inequality gives
\[
\P(S_N\ge(1+\delta)\mu)
\le\exp\left(-\mu\{(1+\delta)\ln(1+\delta)-\delta\}\right).
\]
Notice that for $s\ge0$,
\[
\frac{d}{ds}\left(\ln(1+s)-\frac{2s}{2+s}\right)
=\frac{s^2}{(1+s)(2+s)^2}\ge0.
\]
The expression in parentheses is zero at $s=0$, and hence
\[
(1+\delta)\ln(1+\delta)-\delta
\ge(1+\delta)\frac{2\delta}{2+\delta}-\delta
=\frac{\delta^2}{2+\delta}.
\]
It follows that
\[
\P(S_N-\mu\ge\delta\mu)
\le\exp\left(-\frac{\delta^2\mu}{2+\delta}\right).
\]

For $0<\delta<1$, use the same MGF bound with
$\lambda=\ln(1-\delta)<0$. Since $x\mapsto e^{\lambda x}$ is decreasing,
Markov's inequality gives
\[
\P(S_N\le(1-\delta)\mu)
\le\exp\left(-\mu\{\delta+(1-\delta)\ln(1-\delta)\}\right).
\]
Also,
\[
\delta+(1-\delta)\ln(1-\delta)
=\int_0^\delta-\ln(1-s)\,ds
\ge\int_0^\delta s\,ds=\frac{\delta^2}{2}.
\]
and hence
\[
\P(S_N-\mu\le-\delta\mu)
\le e^{-\delta^2\mu/2}
\le\exp\left(-\frac{\delta^2\mu}{2+\delta}\right).
\]
If $\delta=1$, then
\[
\P(S_N\le0)=\prod_{i=1}^N(1-p_i)\le e^{-\mu}\le e^{-\mu/3}.
\]
If $\delta>1$, the lower-tail event is empty since $S_N\ge0$.
Combining the two tails by the union bound, we obtain
\[
\P(|S_N-\mu|\ge\delta\mu)
\le2\exp\left(-\frac{\delta^2\mu}{2+\delta}\right)
\]
for every $\delta\ge0$, and we are done.
\vspace{1em}

\subsubsection*{Exercise 2.38. Maximum of Gaussians (optional)}
\begin{enumerate}[label=(\alph*),leftmargin=*]
\item Let $g_1,\ldots,g_N$ be $N(0,1)$ random variables for some $N\ge2$,
that are not necessarily independent. Prove that
\[
\E\max_{i\le N}g_i\le\sqrt{2\ln N},
\qquad
\E\max_{i\le N}|g_i|\le\sqrt{2\ln(2N)}.
\]
\item Prove that if $g_1,g_2,\ldots$ are independent $N(0,1)$, then
as $N\to\infty$,
\[
\E\max_{i\le N}g_i=\sqrt{2\ln N}\,(1+o(1)),
\qquad
\E\max_{i\le N}|g_i|=\sqrt{2\ln N}\,(1+o(1)).
\]
\end{enumerate}

\par\vspace{0.5em}\noindent\textbf{Sol.}\par
(a) Denote $M_N=\max_{i\le N}g_i$. For any $\lambda>0$, Jensen's
inequality gives
\[
\begin{aligned}
e^{\lambda\E M_N}
&\le\E e^{\lambda M_N}\\
&\le\sum_{i=1}^N\E e^{\lambda g_i}
=Ne^{\lambda^2/2},
\end{aligned}
\]
and hence
\[
\E M_N\le\frac{\ln N}{\lambda}+\frac{\lambda}{2}.
\]
Taking $\lambda=\sqrt{2\ln N}$, we obtain
$\E M_N\le\sqrt{2\ln N}$.

For $A_N=\max_{i\le N}|g_i|$, notice that
\[
e^{\lambda A_N}
\le\sum_{i=1}^N\bigl(e^{\lambda g_i}+e^{-\lambda g_i}\bigr),
\]
and the same argument gives
\[
\E A_N\le\frac{\ln(2N)}{\lambda}+\frac{\lambda}{2}
\le\sqrt{2\ln(2N)}
\]
when $\lambda=\sqrt{2\ln(2N)}$. Neither argument uses independence.

\noindent (b) It remains to obtain a lower bound. For $g\sim N(0,1)$ and
$t\ge1$, we have
\[
\begin{aligned}
\P(g>t)
&=\frac1{\sqrt{2\pi}}\int_t^\infty e^{-x^2/2}\,dx\\
&\ge\frac1{t\sqrt{2\pi}}
\exp\left(-\frac{(t+1/t)^2}{2}\right)\\
&\ge\frac{c}{t}e^{-t^2/2},
\end{aligned}
\]
where $c=e^{-3/2}/\sqrt{2\pi}>0$. The second line follows by integrating
only over $[t,t+1/t]$.

For any fixed $\epsilon\in(0,1)$, denote
$t_N=(1-\epsilon)\sqrt{2\ln N}$. For $N$ sufficiently large, $t_N\ge1$,
and hence
\[
N\P(g>t_N)
\ge\frac{c}{(1-\epsilon)\sqrt{2\ln N}}\,
N^{1-(1-\epsilon)^2}\longrightarrow\infty.
\]
By independence, we know
\[
\P(M_N\le t_N)
=\bigl(1-\P(g>t_N)\bigr)^N
\le\exp\bigl(-N\P(g>t_N)\bigr)\longrightarrow0.
\]
To pass to expectations, notice that $M_N\ge g_1$ and hence
$M_N^-\le g_1^-$, where $x^-=\max\{-x,0\}$. It follows that
\[
\begin{aligned}
\E M_N
&\ge t_N\P(M_N>t_N)-\E M_N^-\\
&\ge t_N\P(M_N>t_N)-\E g_1^-.
\end{aligned}
\]
Since $\E g_1^-<\infty$, we obtain
\[
\liminf_{N\to\infty}\frac{\E M_N}{\sqrt{2\ln N}}\ge1-\epsilon.
\]
Letting $\epsilon\downarrow0$ and using (a), we know
$\E M_N/\sqrt{2\ln N}\to1$. Finally,
\[
\E M_N\le\E A_N\le\sqrt{2\ln(2N)},
\]
and hence $\E A_N/\sqrt{2\ln N}\to1$ as well.

\newpage
\subsection*{Problem 3. Subgaussian norms and covariance (required)}
In this problem, show that a subgaussian norm bound controls the second
moment, whereas a second-moment bound alone does not imply subgaussianity.
Use the convention
\[
\norm{Z}_{\psi_2}
=\inf\{t>0:\E\exp(Z^2/t^2)\le2\},
\qquad\inf\varnothing=\infty.
\]
This is equivalent, up to absolute constants, to the moment-based
convention used in the lecture notes. Here $\norm{\cdot}_{\mathrm{op}}$
denotes the Euclidean operator norm.

\begin{enumerate}[label=(\arabic*),leftmargin=*]
\item Let $X$ be a real-valued random variable with $\E X=0$.
Prove that $\norm{X}_{\psi_2}\le1$ implies that $\operatorname{Var}(X)$
is bounded by an absolute constant. Give a mean-zero example with
$\operatorname{Var}(X)\le1$ that is not subgaussian.

\par\vspace{0.5em}\noindent\textbf{Sol.}\par
For every $t>1$, the definition of the norm and monotonicity in $t$ give
$\E e^{X^2/t^2}\le2$. Notice that this also implies $\E X^2<\infty$.
By Jensen's inequality,
\[
\exp\left(\frac{\E X^2}{t^2}\right)
\le\E e^{X^2/t^2}\le2,
\]
and hence $\E X^2\le t^2\ln2$. Letting $t\downarrow1$, we obtain
\[
\operatorname{Var}(X)=\E X^2\le\ln2.
\]
This argument does not require the infimum in the definition to be attained.

For the converse, consider $Z$ with density
\[
f_Z(z)=\frac1{\sqrt2}e^{-\sqrt2|z|},\qquad z\in\R.
\]
By symmetry, $\E Z=0$, and
\[
\operatorname{Var}(Z)=\E Z^2
=\sqrt2\int_0^\infty z^2e^{-\sqrt2z}\,dz=1.
\]
However, for every $t>0$,
\[
\E e^{Z^2/t^2}
=\sqrt2\int_0^\infty
\exp\left(\frac{z^2}{t^2}-\sqrt2z\right)\,dz=\infty.
\]
For $z\ge2\sqrt2t^2$ the exponent is at least $z^2/(2t^2)$,
and hence $\norm{Z}_{\psi_2}=\infty$, which gives the required example.
\vspace{1em}

\newpage
\item Let $X\in\R^p$ have $\E X=0$, and let
$S^{p-1}=\{u\in\R^p:\norm{u}_2=1\}$. A random vector is called
subgaussian when $\sup_{u\in S^{p-1}}\norm{\langle X,u\rangle}_{\psi_2}<\infty$.
Prove that
\[
\sup_{u\in S^{p-1}}\norm{\langle X,u\rangle}_{\psi_2}\le1
\]
implies an absolute bound on $\norm{\operatorname{Cov}(X)}_{\mathrm{op}}$,
independent of $p$. Give a mean-zero example that is not a subgaussian
random vector but satisfies $\norm{\operatorname{Cov}(X)}_{\mathrm{op}}\le1$.

\par\vspace{0.5em}\noindent\textbf{Sol.}\par
For every $u\in S^{p-1}$, we know
$\E\langle X,u\rangle=0$ and
$\norm{\langle X,u\rangle}_{\psi_2}\le1$. By (1),
\[
\E\langle X,u\rangle^2\le\ln2.
\]
In particular, taking $u=e_j$ shows that every coordinate has finite
second moment, so the covariance matrix $\Sigma=\E XX^\top$ is defined.
It is symmetric and nonnegative definite, and hence
\[
\norm{\Sigma}_{\mathrm{op}}
=\sup_{u\in S^{p-1}}u^\top\Sigma u
=\sup_{u\in S^{p-1}}\E\langle X,u\rangle^2
\le\ln2.
\]
The constant does not depend on $p$.

For the example, let $Z_1,\ldots,Z_p$ be independent copies of the
variable $Z$ in (1), and denote $X=(Z_1,\ldots,Z_p)^\top$. Then
\[
\E X=0,\qquad\operatorname{Cov}(X)=I_p,
\qquad\norm{\operatorname{Cov}(X)}_{\mathrm{op}}=1.
\]
But $\langle X,e_1\rangle=Z_1$ is not subgaussian, and hence
\[
\sup_{u\in S^{p-1}}\norm{\langle X,u\rangle}_{\psi_2}
\ge\norm{Z_1}_{\psi_2}=\infty.
\]
This gives the required vector example.
\end{enumerate}

\end{document}
